% !PRINTED 324
\weberpage{324}
Irgend eine ternäre Form $n^{\text{ten}}$ Grades $f(x_1, x_2, x_3)$ stellt,
gleich Null gesetzt, eine Curve $n^{\text{ter}}$ Ordnung dar, die mit einer
geraden Linie $n$ Schnittpunkte hat. Diese Curve bezeichnen wir
kurz als die Curve $f$.

Durch die Substitution (1) geht $f$ in eine andere Form desselben Grades in den Variablen $y_1, y_2, y_3$ über:

\weq{3}{f(x_1, x_2, x_3) = \varphi(y_1, y_2, y_3),}

\noindent und $\varphi = 0$ stellt dieselbe Curve dar, wie $f = 0$, bezogen auf
das Coordinatensystem $y_1, y_2, y_3$.

Bilden wir die Ableitungen der Identität (3) nach den
Variablen $y_1, y_2, y_3$, so folgt mit Rücksicht auf (1):

\weq{4}{\begin{aligned}
\varphi'(y_1) &= \alpha_1^{(1)} f'(x_1) + \alpha_2^{(1)} f'(x_2) + \alpha_3^{(1)} f'(x_3) \\
\varphi'(y_2) &= \alpha_1^{(2)} f'(x_1) + \alpha_2^{(2)} f'(x_2) + \alpha_3^{(2)} f'(x_3) \\
\varphi'(y_3) &= \alpha_1^{(3)} f'(x_1) + \alpha_2^{(3)} f'(x_2) + \alpha_3^{(3)} f'(x_3).
\end{aligned}}

\begin{center}
\S.\,86.

\gesperrt{Singuläre Punkte. Wendepunkte. Doppeltangenten.}
\end{center}

Wenn für einen Punkt $x$ die drei Gleichungen

\weq{1}{f'(x_1) = 0, \quad f'(x_2) = 0, \quad f'(x_3) = 0}

\noindent zugleich erfüllt sind, so liegt nach dem \gesperrt{Euler'schen Theorem}

\weq{2}{n f = x_1 f'(x_1) + x_2 f'(x_2) + x_3 f'(x_3)}

\noindent (Bd. I, \S.\,17) dieser Punkt auf der Curve $f$. Die Gleichungen \S.\,85,(4)
zeigen, dass in demselben Punkte auch $\varphi'(y_1), \varphi'(y_2), \varphi'(y_3)$
verschwinden müssen, dass also die durch die Gleichungen (1)
charakterisirte Eigenschaft eines Punktes bei linearer Transformation erhalten bleibt, also, wie wir uns auch ausdrücken, zu
den \gesperrt{invarianten} Eigenschaften gehört.

Nicht auf jeder Curve $n^{\text{ter}}$ Ordnung kommen Punkte vor, die
den Bedingungen (1) genügen, wie z.\,B. die Annahme $f(x_1, x_2, x_3)$
$= x_1^n + x_2^n + x_3^n$ zeigt. Es wird vielmehr, wenn auch nur ein
solcher Punkt existiren soll, eine Gleichung zwischen den Co\"efficienten von $f$ erfüllt sein müssen, die man durch Elimination
von $x_1, x_2, x_3$ aus den drei Gleichungen (1) erhält (Bd. I, \S.\,50).
Man kann diese Relation in der Form darstellen, dass eine
gewisse ganze rationale und homogene Function der sämmtlichen
Co\"efficienten von $f$ gleich Null sein muss, und diese Function
%CONT>>
